\documentclass[11pt,a4paper]{article}
\usepackage[margin=2.6cm]{geometry}
\usepackage{amsmath,amssymb,amsthm,mathtools,booktabs,microtype}
\usepackage{xcolor}
\usepackage[colorlinks,linkcolor=blue!50!black,citecolor=blue!50!black]{hyperref}
\newtheorem{proposition}{Proposition}
\newtheorem{theorem}[proposition]{Theorem}
\newtheorem{lemma}[proposition]{Lemma}
\newtheorem{corollary}[proposition]{Corollary}
\theoremstyle{remark}\newtheorem{remark}[proposition]{Remark}
\newcommand{\dd}{\mathrm{d}}
\newcommand{\E}{\mathbb{E}}
\newcommand{\Var}{\operatorname{Var}}
\newcommand{\Cov}{\operatorname{Cov}}
\newcommand{\Ain}{A_{\infty}}
\newcommand{\Bin}{B_{\infty}}
\newcommand{\nustat}{\nu_{\mathrm{stat}}}
\title{The stationary-start maintenance model\\ and its exact delay-period predictions}
\author{Working note}
\date{\today}
\begin{document}\maketitle

\begin{abstract}\noindent
One change to the maintenance model --- initialising the Ornstein--Uhlenbeck
maintenance noise at its stationary law rather than at zero --- makes the transient
irrelevant to the delay-period analysis window, and every prediction the paper tests
then follows from the \emph{stationary} Kac--Rice rate by a short corollary. The note
is organised accordingly: Sections~\ref{sec:model}--\ref{sec:neglect} establish that
the transient can be neglected, and state precisely under what conditions;
Sections~\ref{sec:stat}--\ref{sec:pred} work with the stationary rate alone and derive
the predictions. Every statement has been verified numerically
(Section~\ref{sec:checks}).
\end{abstract}

\tableofcontents

\section{The model}\label{sec:model}

Everything before stimulus offset is unchanged. Maintenance begins at
$\tau=t_{\mathrm{offset}}+L_{\mathrm{motion}}$, at which time the gate-facing signal
$y$ equals the maintained target $m$. During maintenance,
\begin{align}
\dd\varepsilon(t)&=-\kappa\,\varepsilon(t)\,\dd t+\sigma_\varepsilon\,\dd W_\varepsilon(t),
\label{eq:ou}\\
\dd y(t)&=\lambda\,[\,m+\varepsilon(t)-y(t)\,]\,\dd t ,
\label{eq:filt}
\end{align}
and $H:=y-m$ obeys $\dd H=\lambda(\varepsilon-H)\,\dd t$.

\paragraph{Initial conditions.}
$H(\tau)=0$ is \emph{forced}: the filter output is continuous and $m$ is defined as
the value of $y$ at $\tau$. It is not a modelling choice. By contrast
$\varepsilon(\tau)$ is free; the present model sets it to $0$. We propose instead
\begin{equation}
\varepsilon(\tau)\sim\mathrm{N}\!\left(0,\ \sigma_\varepsilon^2/2\kappa\right)
\quad\text{independent of }W_\varepsilon
\tag{\textsc{ss}}\label{eq:warm}
\end{equation}
--- the stationary law of \eqref{eq:ou}. Call \eqref{eq:warm} the \emph{stationary
start} and $\varepsilon(\tau)=0$ the \emph{cold start}.

\paragraph{Why.}
$\varepsilon$ represents fluctuations intrinsic to holding a value. Nothing makes
such a process sit exactly at its mean at the instant maintenance begins, and no
measurement constrains it. As Corollary~\ref{cor:rate} shows, the cold start is the
sole source of a slow ($1/\kappa$) transient in a system whose only other timescale
is the fast filter $1/\lambda$; under \eqref{eq:warm} the model's delay-period
predictions become exact rather than asymptotic.

\begin{remark}
\eqref{eq:warm} removes the mechanism that gives Model~1 its only delay-period ramp.
That is not a way of disarming Model~1: Section~\ref{sec:neglect} bounds what
\emph{any} residual non-stationarity can buy it, and the bound is far below the data.
\end{remark}

\paragraph{Notation.}
$\rho:=t-\tau$ is elapsed maintenance time; $\beta/\delta$ is the speed of temporal
preparation, with $\beta=0$ under Model~1; $d(t)=g-|m|-(\beta/\delta)(t-t_{\mathrm{offset}})^{+}$
is the distance to the gate; $\varphi,\Phi$ are the standard normal density and
c.d.f., and $g(u):=u\Phi(u)+\varphi(u)$, so $g'=\Phi$ and $g''=\varphi$. Write
$s^2:=\sigma_\varepsilon^2/2\kappa$ and
\begin{equation}
\Ain:=\frac{\lambda s^2}{\lambda+\kappa},\qquad
\Bin:=\frac{\lambda^{2}\sigma_\varepsilon^{2}}{2(\lambda+\kappa)},\qquad
\omega:=\sqrt{\Bin/\Ain}=\sqrt{\lambda\kappa},\qquad
a_\infty:=\frac{d}{\sqrt{\Ain}} .
\label{eq:stat}
\end{equation}

\section{Exact second-order moments}\label{sec:moments}

\begin{proposition}[exact moments]\label{prop:moments}
Under \eqref{eq:ou}--\eqref{eq:warm} with $H(\tau)=0$, put
\begin{equation}
X(\rho):=\frac{(\lambda+\kappa)\bigl(e^{-2\lambda\rho}-e^{-(\lambda+\kappa)\rho}\bigr)}{\lambda-\kappa}
\qquad\bigl(=-2\lambda\rho\,e^{-2\lambda\rho}\ \text{ when }\lambda=\kappa\bigr).
\label{eq:Xdef}
\end{equation}
Then for every $\rho>0$,
\begin{equation}
\boxed{\ \
\frac{\sigma_H^2}{\Ain}=1+X-e^{-(\lambda+\kappa)\rho},\qquad
\frac{c_H}{\Ain}=-\lambda X,\qquad
\frac{\sigma_{\dot H}^2}{\Bin}=1+\frac{\lambda}{\kappa}\Bigl(X+e^{-(\lambda+\kappa)\rho}\Bigr).
\ \ }
\label{eq:moments}
\end{equation}
In particular $\sigma_H^2(0)=0$, $c_H(0)=0$,
$\sigma_{\dot H}^2(0)=\frac{\lambda+\kappa}{\kappa}\Bin$, and all three converge to
\eqref{eq:stat} with $c_H(\infty)=0$.
\end{proposition}

\begin{proof}
$(\varepsilon,H)$ is a two-dimensional linear diffusion with
$F=\bigl(\begin{smallmatrix}-\kappa&0\\ \lambda&-\lambda\end{smallmatrix}\bigr)$ and
$Q=\mathrm{diag}(\sigma_\varepsilon^2,0)$, so its covariance solves
$\dot P=FP+PF^{\!\top}+Q$ with $P(0)=\mathrm{diag}(s^2,0)$. The system is triangular:
\begin{align}
\dot P_{\varepsilon\varepsilon}&=-2\kappa P_{\varepsilon\varepsilon}+\sigma_\varepsilon^2,
&&P_{\varepsilon\varepsilon}(0)=s^2,\label{eq:l1}\\
\dot P_{\varepsilon H}&=-(\kappa+\lambda)P_{\varepsilon H}+\lambda P_{\varepsilon\varepsilon},
&&P_{\varepsilon H}(0)=0,\label{eq:l2}\\
\dot P_{HH}&=2\lambda(P_{\varepsilon H}-P_{HH}),&&P_{HH}(0)=0.\label{eq:l3}
\end{align}
By \eqref{eq:l1}, $\dot P_{\varepsilon\varepsilon}(0)=-2\kappa s^2+\sigma_\varepsilon^2=0$,
so $P_{\varepsilon\varepsilon}\equiv s^2$: \emph{the stationary start freezes the first
component exactly.} Then \eqref{eq:l2} gives
$P_{\varepsilon H}=\Ain\bigl(1-e^{-(\lambda+\kappa)\rho}\bigr)$, and \eqref{eq:l3},
integrated from $0$,
\[
P_{HH}=2\lambda\!\int_0^{\rho}\!e^{-2\lambda(\rho-v)}P_{\varepsilon H}(v)\dd v
=\Ain\left[1-e^{-2\lambda\rho}-2\lambda\frac{e^{-(\lambda+\kappa)\rho}-e^{-2\lambda\rho}}{\lambda-\kappa}\right],
\]
and writing $-2\lambda=-(\lambda+\kappa)-(\lambda-\kappa)$ in the numerator gives
$P_{HH}/\Ain-1=X-e^{-(\lambda+\kappa)\rho}$. Since $\dot H=\lambda(\varepsilon-H)$,
$c_H=\lambda(P_{\varepsilon H}-P_{HH})=-\lambda\Ain X$ and
$\sigma_{\dot H}^2=\lambda^{2}(P_{\varepsilon\varepsilon}-2P_{\varepsilon H}+P_{HH})
=\Bin+\lambda^{2}\Ain\bigl(X+e^{-(\lambda+\kappa)\rho}\bigr)$, using
$\lambda^2(s^2-\Ain)=\Bin$ and $\lambda^{2}\Ain/\Bin=\lambda/\kappa$. The $\lambda=\kappa$
forms follow by l'H\^opital.
\end{proof}

\begin{remark}
$\sigma_H$ rises from $0$; $\sigma_{\dot H}$ \emph{falls} from
$\frac{\lambda+\kappa}{\kappa}\Bin$, because $\dot H(\tau)=\lambda\varepsilon(\tau)$
has the noise at full amplitude while $H$ is still at zero. The position variance
fills in while the velocity variance settles down.
\end{remark}

\section{How fast the transient dies}\label{sec:relax}

\begin{lemma}\label{lem:mvt}
For $A,B>0$ and $\rho>0$,
$\bigl|e^{-A\rho}-e^{-B\rho}\bigr|\le|A-B|\,\rho\,e^{-\min(A,B)\rho}$.
\end{lemma}
\begin{proof}
The mean value theorem applied to $x\mapsto e^{-x\rho}$ on the interval with
endpoints $A,B$ gives $|A-B|\rho e^{-\xi\rho}$ for some $\xi$ between them, and
$e^{-\xi\rho}\le e^{-\min(A,B)\rho}$.
\end{proof}

\begin{proposition}[transient bounds]\label{prop:bounds}
Set $\mu:=\lambda+\min(\lambda,\kappa)$, $E_1:=e^{-(\lambda+\kappa)\rho}$,
$E_2:=e^{-2\lambda\rho}$, and
\begin{equation}
\mathcal{D}(\rho):=\min\left\{\rho\,e^{-\mu\rho},\ \frac{E_1+E_2}{|\lambda-\kappa|}\right\}
\qquad(\text{first entry only, when }\lambda=\kappa).
\label{eq:Dcal}
\end{equation}
Then, for all $\lambda,\kappa,\rho>0$,
\begin{equation}
\bar\Delta:=(\lambda+\kappa)\mathcal{D}+E_1,\qquad
\bar\gamma:=\lambda(\lambda+\kappa)\mathcal{D},\qquad
\bar e:=\min\left\{\frac{\lambda}{\kappa}\bar\Delta,\
\frac{\lambda\bigl[(\lambda+\kappa)E_2+2\kappa E_1\bigr]}{\kappa|\lambda-\kappa|}\right\}
\label{eq:Dbar}
\end{equation}
bound the three deviations:
\begin{equation}
\left|\frac{\sigma_H^2}{\Ain}-1\right|\le\bar\Delta,\qquad
0\le\frac{c_H}{\Ain}\le\bar\gamma,\qquad
\left|\frac{\sigma_{\dot H}^2}{\Bin}-1\right|\le\bar e .
\label{eq:threebounds}
\end{equation}
All three are decreasing on $\rho\ge1$ when $\lambda\ge1$.
\end{proposition}
\begin{proof}
Everything reduces to bounding $|E_1-E_2|/|\lambda-\kappa|$, since by
Proposition~\ref{prop:moments} $X=(\lambda+\kappa)(E_2-E_1)/(\lambda-\kappa)$,
$\sigma_H^2/\Ain-1=X-E_1$, $c_H/\Ain=-\lambda X$, and
$\sigma_{\dot H}^2/\Bin-1=(\lambda/\kappa)\bigl[(\lambda+\kappa)E_2-2\kappa E_1\bigr]/(\lambda-\kappa)$,
this last by adding and subtracting $(\lambda+\kappa)E_1$ in the numerator.
Two bounds are available and we take the smaller. \emph{(a)} By
Lemma~\ref{lem:mvt} with $A=2\lambda$, $B=\lambda+\kappa$ --- so $|A-B|=|\lambda-\kappa|$
and $\min(A,B)=\mu$ --- $|E_1-E_2|/|\lambda-\kappa|\le\rho e^{-\mu\rho}$.
\emph{(b)} Trivially $|E_1-E_2|\le E_1+E_2$. Bound (a) is the sharp one when
$\lambda\approx\kappa$, where (b) is singular; (b) is the sharp one when the two
rates are far apart, where (a)'s worst-case exponent $\mu$ is pessimistic. The sign
$X\le0$ follows by inspection of \eqref{eq:Xdef}, giving $c_H\ge0$. For the third
bound, the first entry of $\bar e$ comes from the $X$ form and the second from
bounding the two exponentials separately. Monotonicity: for the first entry,
$\frac{\dd}{\dd\rho}(1+c\rho)e^{-\mu\rho}=e^{-\mu\rho}(c-\mu-\mu c\rho)$ with
$c=\lambda+\kappa\ge\mu$, negative for $\rho>(c-\mu)/(\mu c)$, and
$(c-\mu)/(\mu c)\le1/(2\lambda)\le1/2$ when $\lambda\ge1$; the second entry is a sum
of decaying exponentials.
\end{proof}

\begin{corollary}[relaxation is at least as fast as $\lambda$]\label{cor:rate}
\emph{From Propositions~\ref{prop:moments} and~\ref{prop:bounds}.}
All three moments converge to \eqref{eq:stat} at rate $\mu\ge\lambda$, uniformly in
$\kappa$. The operator $P\mapsto FP+PF^{\!\top}$ has eigenvalues
$\{-2\kappa,-(\lambda+\kappa),-2\lambda\}$, so a general initial condition relaxes as
a combination of $e^{-2\kappa\rho}$ (the $\varepsilon$--$\varepsilon$ block),
$e^{-(\lambda+\kappa)\rho}$ (the cross term) and $e^{-2\lambda\rho}$ (the
$H$--$H$ block). Under \eqref{eq:warm} the first has zero coefficient, because
$P_{\varepsilon\varepsilon}$ is constant; the slower of the remaining two is
$e^{-\mu\rho}$, which equals $e^{-2\lambda\rho}$ only when $\kappa\ge\lambda$.
\end{corollary}

\begin{remark}[the $\kappa$ floor]\label{rem:kappa}
The first two bounds in \eqref{eq:threebounds} are uniform in $\kappa$; the third
carries $\lambda/\kappa$ and is not. What it needs is weak: at $\rho=1$ the deviation
is $(\lambda/\kappa)e^{-2\lambda}$ to leading order, so relative accuracy $\epsilon$
needs $\kappa\gtrsim\lambda e^{-2\lambda}/\epsilon$, which at $\lambda=10$,
$\epsilon=10^{-2}$ is $\kappa\ge2\times10^{-6}$. Read through the correlation time:
$\Cov(\varepsilon(t),\varepsilon(t+r))\propto e^{-\kappa|r|}$ at stationarity, so
$1/\kappa$ is the lag at which the maintenance noise decorrelates to $1/e$. A process
with $1/\kappa$ of minutes is constant over the $3$~s delay and over dozens of
consecutive trials --- a session-level drift, for which the model has no term. A floor
$\kappa\ge0.05\,\mathrm{s}^{-1}$ ($1/\kappa\le20$~s, about seven delay periods) is
therefore a statement about what the term means, not a range imported from simulation.
\end{remark}

\section{Neglecting the transient}\label{sec:neglect}

Let $\nustat$ denote the rate computed from the stationary moments \eqref{eq:stat},
i.e.\ with $\sigma_H^2=\Ain$, $\sigma_{\dot H}^2=\Bin$, $c_H=0$. This section
establishes the single result the rest of the note rests on.

\begin{theorem}[the stationary approximation]\label{thm:main}
Assume
\begin{equation}
\text{\emph{(i)}}\ \lambda\ge5\,\mathrm{s}^{-1},\qquad
\text{\emph{(ii)}}\ \kappa\ge0.2\,\mathrm{s}^{-1},\qquad
\text{\emph{(iii)}}\ a_\infty\le5 .
\label{eq:conds}
\end{equation}
Then, throughout the analysis window $\rho\in[1,3]$~s and for both models:
\begin{enumerate}
\item[\emph{(a)}] $\nu(\rho)=\nustat(\rho)\,e^{\theta(\rho)}$ with
$|\theta(\rho)|\le\mathcal{E}(\rho)$ and $\mathcal{E}(1)\le0.36$ --- the rate is
exact to within a factor $1.43$, and to within $1.0004$ by the end of the window;
\item[\emph{(b)}] the least-squares slope and quadratic coefficient of $\ln\nu$ on the
$20$ bins of the window differ from those of $\ln\nustat$ by at most
\[
\bigl|c_1-c_1^{\mathrm{stat}}\bigr|\le0.081\ \mathrm{s}^{-1},
\qquad
\bigl|c_2-c_2^{\mathrm{stat}}\bigr|\le0.145\ \mathrm{s}^{-2};
\]
\item[\emph{(c)}] under $\lambda\ge6$ these become $0.027\,\mathrm{s}^{-1}$ and
$0.049\,\mathrm{s}^{-2}$; under $\lambda\ge8$, $0.0030$ and $0.0058$.
\end{enumerate}
\end{theorem}

\noindent The observed log-rate slope is $+0.867\,\mathrm{s}^{-1}$, so (b) is a
$9\%$ error bar on it and (c) a $3\%$ one. \emph{Everything from
Section~\ref{sec:stat} on assumes \eqref{eq:conds} and works with $\nustat$, carrying
(b) as the error bar.} Sections~\ref{sec:conds} and~\ref{sec:proof} justify the
hypotheses and prove the theorem.

\subsection{The three conditions}\label{sec:conds}

Each can be argued without reference to the measured rate.

\begin{description}
\item[(i) $\lambda\ge5\,\mathrm{s}^{-1}$, i.e.\ $1/\lambda\le200$~ms.]
$1/\lambda$ is the time constant of the low-pass filter carrying motor preparation to
the gate. Dendauw et al.\ estimate $\lambda\approx20\,\mathrm{s}^{-1}$ ($50$~ms) in an
immediate-report task, so \eqref{eq:conds} is four times more permissive than the
published value. $\lambda$ is in principle identifiable from these data, since it
sets how fast the rate follows the stimulus; that check has not been carried out.
\item[(ii) $\kappa\ge0.2\,\mathrm{s}^{-1}$, i.e.\ $1/\kappa\le5$~s.]
$1/\kappa$ is the correlation time of the maintenance noise:
$\Cov(\varepsilon(t),\varepsilon(t+r))\propto e^{-\kappa|r|}$ at stationarity, so it
is the lag at which the fluctuation has decorrelated to $1/e$. Requiring it to be at
most one trial says that what is maintained fluctuates \emph{within} a trial; a longer
correlation time describes a drift carried across trials, for which the model has no
term and which the randomised design averages out rather than expresses as a
delay-period ramp. (i) and (ii) are substitutes: $\lambda\to\lambda+1$ buys about what
$\kappa\to2\kappa$ buys, so only one needs to be argued hard.
\item[(iii) $a_\infty\le5$.] Not fitted. By \eqref{eq:m1flat} the expected number of
partial activations per $2$~s delay window is $(\omega/\pi)e^{-a^{2}/2}$; at the
largest $\omega$ in play this is one per $19$ trials at $a=3.3$, one per $246$ at
$a=4$, and \emph{one per $22\,000$ at $a=5$}. So (iii) says only that the model must
produce partial activations at all --- the experiment comprises of order $10^{4}$
trials --- not that it must reproduce the observed rate. Section~\ref{sec:afromnu}
gives the sharper, data-based $a\le3.3$, which is not used.
\end{description}

\subsection{Proof of Theorem~\ref{thm:main}}\label{sec:proof}

\begin{proposition}[pointwise error]\label{prop:point}
Fix $\rho\ge1$, evaluate $\bar\Delta,\bar\gamma,\bar e$ of \eqref{eq:Dbar} at $\rho$,
and write $u_s:=\beta/(\delta\sqrt{\Bin})$, $c_g:=\Phi(0)/g(0)=1.2533$ and
\begin{equation}
\eta:=\bar e+\frac{\bar\gamma^{2}}{\omega^{2}(1-\bar\Delta)} .
\label{eq:eta}
\end{equation}
If $\bar\Delta<1$ and $\eta<1$ then, for both models,
\begin{equation}
\bigl|\ln\nu(\rho)-\ln\nustat(\rho)\bigr|\ \le\ \mathcal{E}(\rho):=
\frac{(a_\infty^{2}+1)\bar\Delta}{2(1-\bar\Delta)}
+\frac{\eta}{2(1-\eta)}
+c_g\left[\frac{\bar\gamma\,a_\infty}{\omega(1-\bar\Delta)\sqrt{1-\eta}}
+u_s\!\left(\frac{1}{\sqrt{1-\eta}}-1\right)\right].
\label{eq:point}
\end{equation}
\end{proposition}

\begin{proof}
Write $\Delta:=\sigma_H^2/\Ain-1$, $\gamma:=c_H/\Ain$, $E:=\sigma_{\dot H}^2/\Bin-1$,
so $|\Delta|\le\bar\Delta$, $0\le\gamma\le\bar\gamma$, $|E|\le\bar e$ by
Proposition~\ref{prop:bounds}. From \eqref{eq:nu},
\[
\ln\nu-\ln\nustat=-\tfrac12\bigl(a^{2}-a_\infty^{2}\bigr)
-\tfrac12\ln\frac{\sigma_H^{2}}{\Ain}+\ln\frac{G}{G_s},
\]
with $G=\sqrt S\,g(u)$ and $G_s=\sqrt{\Bin}\,g(u_s)$. Throughout,
$|\ln(1+x)|\le|x|/(1-|x|)$ for $|x|<1$.

\emph{(1)} $a^{2}=a_\infty^{2}/(1+\Delta)$, so
$\bigl|a^{2}-a_\infty^{2}\bigr|\le a_\infty^{2}\bar\Delta/(1-\bar\Delta)$; halving and
adding $\tfrac12|\ln(1+\Delta)|\le\bar\Delta/(2(1-\bar\Delta))$ gives the first term of
\eqref{eq:point}.

\emph{(2)} $S=\Bin\bigl[(1+E)-\gamma^{2}/(\omega^{2}(1+\Delta))\bigr]$ since
$\Ain/\Bin=\omega^{-2}$, so $|S/\Bin-1|\le\eta$, hence $S\ge\Bin(1-\eta)>0$ and
$\tfrac12|\ln(S/\Bin)|\le\eta/(2(1-\eta))$: the second term.

\emph{(3)} $(\ln g)'=\Phi/g$ and $(\ln g)''=-D/g^{2}<0$, where $D:=\Phi^{2}-\varphi g$
satisfies $D'=\varphi(\Phi+ug)>0$ and $D(-\infty)=0$, hence $D>0$; so $\Phi/g$ is
decreasing and $\sup_{u\ge0}(\ln g)'=\Phi(0)/g(0)=c_g$. Since
$M=\beta/\delta+\gamma d/(1+\Delta)\ge\beta/\delta\ge0$ we have $u,u_s\ge0$, so
$|\ln g(u)-\ln g(u_s)|\le c_g|u-u_s|$. Splitting
$u-u_s=\frac{M-\beta/\delta}{\sqrt S}+\frac{\beta}{\delta}\bigl(\frac1{\sqrt S}-\frac1{\sqrt{\Bin}}\bigr)$,
the first piece is at most
$\bar\gamma a_\infty/\bigl(\omega(1-\bar\Delta)\sqrt{1-\eta}\bigr)$ using
$M-\beta/\delta=\gamma d/(1+\Delta)$, $d=a_\infty\sqrt{\Ain}$ and
$\sqrt S\ge\sqrt{\Bin(1-\eta)}$; the second equals
$u_s|(1+x)^{-1/2}-1|$ with $|x|\le\eta$, hence at most
$u_s(1/\sqrt{1-\eta}-1)$. Multiplying by $c_g$ gives the third term.
\end{proof}

\begin{proposition}[transfer to the measured coefficients]\label{prop:transfer}
On the $20$ bin centres $\rho_1,\dots,\rho_{20}$ of the window, the least-squares
slope and quadratic coefficient are linear functionals
$c_j=\sum_i w^{(j)}_i f_i$ of the profile, with
$\sum_i|w^{(\mathrm{slope})}_i|=1.504$ and $\sum_i|w^{(2)}_i|=2.916$. Hence
\begin{equation}
\bigl|c_j-c_j^{\mathrm{stat}}\bigr|\ \le\ \sum_i\bigl|w^{(j)}_i\bigr|\,\mathcal{E}(\rho_i)
\ \le\ \Bigl(\textstyle\sum_i\bigl|w^{(j)}_i\bigr|\Bigr)\,\mathcal{E}(1).
\label{eq:transfer}
\end{equation}
The first form is the one used: $\mathcal{E}$ decays like $e^{-\mu\rho}$ across the
window, so it is $7$ to $10$ times smaller than the uniform second form.
\end{proposition}
\begin{proof}
Immediate from Proposition~\ref{prop:point} and $\mathcal{E}$ decreasing
(Proposition~\ref{prop:bounds}).
\end{proof}

\begin{proof}[Proof of Theorem~\ref{thm:main}]
Under \eqref{eq:conds}, $\mathcal{E}$ of \eqref{eq:point} is increasing in
$a_\infty$ and, at $\rho=1$, maximised over $\kappa\ge0.2$ at $\kappa=0.2$; evaluating
gives $\mathcal{E}(1)=0.353$ at $\lambda=5$, $0.133$ at $\lambda=6$ and $0.019$ at
$\lambda=8$, which is (a). Feeding $\mathcal{E}(\rho_i)$ into \eqref{eq:transfer}
gives (b) and (c):
\begin{center}
\begin{tabular}{rrrrr}
\toprule
$\lambda\ge$ & $1/\lambda$ & $\mathcal{E}(1)$ & slope error & curvature error\\
\midrule
$4$ & $250$ ms & $0.96$  & $0.253$ & $0.446$\\
$5$ & $200$ ms & $0.353$ & $0.081$ & $0.145$\\
$6$ & $167$ ms & $0.133$ & $0.027$ & $0.049$\\
$8$ & $125$ ms & $0.019$ & $0.0030$ & $0.0058$\\
\bottomrule
\end{tabular}
\end{center}
\end{proof}

\begin{remark}[why the transfer step is needed]\label{rem:whyneeded}
Proposition~\ref{prop:bounds} bounds the \emph{moments}; the observables are a slope
and a curvature, that is, \emph{derivatives}, and a bound on a function does not bound
its derivative. Proposition~\ref{prop:point} bridges the gap by bounding $\ln\nu$
itself, pointwise, after which \eqref{eq:transfer} is elementary because $c_1$ and
$c_2$ are bounded linear functionals of the profile.
\end{remark}

\section{The stationary rate}\label{sec:stat}

Let $Z:=s_{\mathrm{resp}}H-d$, so that partial activations are upcrossings of zero by
$Z$ and $\dot Z=\dot H+\beta/\delta$. The Kac--Rice formula for a differentiable
Gaussian process gives
\begin{equation}
\nu(t)=p_{Z(t)}(0)\,\E\bigl[\dot Z\mathbf{1}_{\{\dot Z>0\}}\mid Z=0\bigr]
=\frac{\varphi(a)}{\sigma_H}\,G,
\qquad
G=M\Phi(u)+\sqrt S\varphi(u)=\sqrt S\,g(u),
\label{eq:nu}
\end{equation}
with $a=d/\sigma_H$, $M=\beta/\delta+(c_H/\sigma_H^{2})d$,
$S=\sigma_{\dot H}^{2}-c_H^{2}/\sigma_H^{2}$ and $u=M/\sqrt S$.

\begin{proposition}[the stationary rate]\label{prop:statrate}
Substituting the stationary moments \eqref{eq:stat} ($c_H=0$) into \eqref{eq:nu},
\begin{equation}
\nustat=\varphi(a)\left[\frac{\beta}{\delta\sigma_H}\Phi\!\left(\frac{\beta}{\delta\sigma_{\dot H}}\right)
+\omega\,\varphi\!\left(\frac{\beta}{\delta\sigma_{\dot H}}\right)\right],
\label{eq:nuasymp}
\end{equation}
whose bracket does not depend on $d$. Consequently:
\begin{enumerate}
\item[(i)] \textbf{Model 2} ($\beta>0$): $a(\rho)=a(0)-\beta\rho/(\delta\sigma_H)$ is
affine, so exactly
\begin{equation}
\frac{\dd\ln\nustat}{\dd\rho}=\frac{\beta}{\delta\sigma_H}\,a(\rho),
\qquad
\frac{\dd^{2}\ln\nustat}{\dd\rho^{2}}=-\left(\frac{\beta}{\delta\sigma_H}\right)^{2},
\label{eq:slope}
\end{equation}
a downward parabola, the same one in all six conditions.
\item[(ii)] \textbf{Model 1} ($\beta=0$): $M=0$, $u=0$, $a$ constant, and
\begin{equation}
\nustat=\frac{\omega}{2\pi}\,e^{-a^{2}/2},\qquad
\frac{\dd\ln\nustat}{\dd\rho}\equiv0,\qquad
\frac{\dd^{2}\ln\nustat}{\dd\rho^{2}}\equiv0 .
\label{eq:m1flat}
\end{equation}
\end{enumerate}
\end{proposition}
\begin{proof}
Set $c_H=0$ in \eqref{eq:nu}; then $\ln\nustat=\mathrm{const}-a^{2}/2$ and
\eqref{eq:slope} follows from $a$ affine. For (ii) set $\beta=0$: $M=0$,
$g(0)=1/\sqrt{2\pi}$, $G=\sqrt{\Bin}/\sqrt{2\pi}=\omega\sigma_H/\sqrt{2\pi}$, so
$\nustat=\varphi(a)\omega/\sqrt{2\pi}=\omega e^{-a^{2}/2}/2\pi$, constant in $\rho$.
\end{proof}

\begin{proposition}[level orders by distance]\label{prop:level}
\emph{From \eqref{eq:nuasymp}.}
Under both models $\nustat=\varphi(a)\times(\text{a factor independent of }d)$, so
$\nustat$ is strictly decreasing in $a$, with no condition whatever. Under Model~2 the
slope $\beta a/(\delta\sigma_H)$ of \eqref{eq:slope} is strictly \emph{increasing} in
$a$. Hence level and slope order oppositely across conditions.
\end{proposition}

\begin{remark}[before the window]\label{rem:prewindow}
The same monotonicity holds during the transient, though it then needs a condition:
from \eqref{eq:nu}, $\partial\ln\nu/\partial d=-a/\sigma_H+c_H\Phi(u)/(\sigma_H^{2}G)$
is negative iff $dG>c_H\Phi(u)$, and $G>M\Phi(u)$ with
$dM=(\beta/\delta)d+c_Ha^{2}\ge c_H$ whenever $a\ge1$. Not needed for anything below.
\end{remark}

\section{The standardised distance is measured}\label{sec:afromnu}

Inverting \eqref{eq:m1flat} --- Model~1's own stationary rate formula, so using it to
bound Model~1 is legitimate ---
\begin{equation}
a=\sqrt{2\ln\!\left(\frac{\omega}{2\pi\nu}\right)} .
\label{eq:afromnu}
\end{equation}
(C3) uses only the weak, data-free reading of this identity --- $a\le5$, on pain of
the model producing fewer than one partial activation in twenty thousand trials.
Two sharper readings are available if one is willing to use the measured rate, and
neither is needed above. First, with the observed delay-period rates
$0.030$--$0.135\,\mathrm{s}^{-1}$ and $(\lambda,\kappa)$ anywhere in the sampled
range, \eqref{eq:afromnu} gives $a\le3.3$, which tightens every ceiling in
Section~\ref{sec:conds} by about a third. Second, \eqref{eq:afromnu} is the
rate--ramp trade-off in its sharpest form: the rate \emph{measures} the standardised
distance, so the ramp Model~1 could produce and the rate it must produce are not
free of one another.

\section{Predictions}\label{sec:pred}

The delay-period observations are: \textbf{(L1)} the level rises with duration;
\textbf{(L2)} the level rises with coherence; \textbf{(L3)} the coherence effect on
the level is larger after the long stimulus; \textbf{(S1)} the log-rate slope is
$+0.867\,\mathrm{s}^{-1}$; \textbf{(S2)} it is shallower after the long stimulus,
$0.896\to0.723$; \textbf{(S3)} shallower at high coherence, $0.871\to0.726$.
Throughout we assume \eqref{eq:conds} and use $\nustat$, with the error bars of
Theorem~\ref{thm:main}(b): $\pm0.081\,\mathrm{s}^{-1}$ on a slope at $\lambda\ge5$,
$\pm0.027$ at $\lambda\ge6$.

\begin{corollary}[L1, L2: the level]\label{cor:P1}
\emph{From Proposition~\ref{prop:level}, plus the saturation argument of the main
text, which concerns the distribution of $m$ and is untouched by \eqref{eq:warm}.}
Duration and coherence each raise the fraction of trials whose maintained target has
reached the cap $\theta$, hence lower $d$ in the sense of first-order stochastic
dominance; since $\nustat$ is decreasing in $d$ at every $\rho$, and averaging over
trials preserves the ordering, both raise the level. Predicted by \emph{both} models:
\emph{not discriminating}.
\end{corollary}

\begin{corollary}[L3: the interaction]\label{cor:P1b}
\emph{From Proposition~\ref{prop:level} and the same saturation argument.}
As the capped fraction tends to one the coherence levels converge, and duration
raises that fraction at every coherence, so the coherence effect on the level must
\emph{shrink} with duration. Predicted by both models, and \emph{contradicted} by the
data. Not discriminating; a genuine disagreement.
\end{corollary}

\begin{corollary}[S1: the slope]\label{cor:P2}
\emph{From Proposition~\ref{prop:statrate}(ii) and Theorem~\ref{thm:main}(b).}
Model~1 predicts a slope of exactly $0$, hence a measured slope of at most
$0.081\,\mathrm{s}^{-1}$ (at $\lambda\ge5$; $0.027$ at $\lambda\ge6$). Model~2
predicts $\beta a(\rho)/(\delta\sigma_H)>0$, its only source of a ramp.
\emph{Discriminating: the observed $+0.867\,\mathrm{s}^{-1}$ exceeds Model~1's
ceiling by a factor of $10.8$ at $\lambda\ge5$ and $33$ at $\lambda\ge6$.}
\end{corollary}

\begin{corollary}[S2, S3: the slope across conditions]\label{cor:P3}
\emph{From Propositions~\ref{prop:statrate},~\ref{prop:level} and~\ref{cor:P1},
with Theorem~\ref{thm:main}(b).}
Model~1 predicts a slope of zero in \emph{every} condition, hence a condition
difference of at most twice the slope error: $0.161\,\mathrm{s}^{-1}$ at
$\lambda\ge5$, $0.053$ at $\lambda\ge6$. Model~2 predicts a difference of definite
sign: the slope increases in $a$ while the level decreases in $a$, so any
manipulation that raises the level lowers the slope, and by Corollary~\ref{cor:P1}
duration and coherence both raise the level. \emph{The observed differences $0.173$
and $0.145$ carry the sign Model~2 requires; at $\lambda\ge6$ they exceed Model~1's
ceiling by $3.3$ and $2.7$, at $\lambda\ge5$ only marginally. This is the leg that
pays for a low $\lambda$; S1 does not.}
\end{corollary}

\begin{remark}[on the curvature]\label{rem:curv}
Under the cold start the curvature of the log rate discriminated, through the decay
that Model~1's variance filling forces. Under \eqref{eq:warm} there is no decay to
force: both models predict a near-zero curvature (Model~1 exactly $0$, Model~2
$-\frac12(\beta/\delta\sigma_H)^{2}$, which is $\ge-0.034$ at the observed slope with
$a\le3.3$), and the observed $-0.0405\pm0.0416$ is consistent with both. The
curvature therefore carries nothing here and is dropped from the list of tested
observations; the discriminating role passes to S2 and S3.
\end{remark}

\section{Summary}

Under \eqref{eq:conds} --- $\lambda\ge5$, $\kappa\ge0.2$, $a_\infty\le5$ --- with the
error bars of Theorem~\ref{thm:main}:
\begin{center}
\begin{tabular}{lllc}
\toprule
& Model 1 & Model 2 & verdict\\
\midrule
L1, L2 level & rises & rises & shared\\
L3 interaction & shrinks & shrinks & shared; both fail\\
S1 slope & $0\pm0.081$ & $\beta a/\delta\sigma_H>0$ & \textbf{discriminating}, $10.8\times$\\
S2, S3 by condition & $0\pm0.161$ & sign forced & sign; magnitude at $\lambda\ge6$\\
\bottomrule
\end{tabular}
\end{center}
At $\lambda\ge6$ the ceilings are $0.027$ and $0.053$, and all three of S1--S3
discriminate by magnitude ($33\times$, $3.3\times$, $2.7\times$).

\section{Numerical verification}\label{sec:checks}

\begin{center}
\begin{tabular}{lll}
\toprule
Claim & Test & Result\\
\midrule
Prop.~\ref{prop:moments} & vs.\ numerical integration of $\dot P=FP+PF^{\!\top}+Q$ & max rel.\ err.\ $8.9\times10^{-12}$\\
Prop.~\ref{prop:moments} & the three identities in terms of $X$ & agree to $1.5\times10^{-13}$\\
Prop.~\ref{prop:bounds} & $0\le-X\le(\lambda+\kappa)\rho e^{-\mu\rho}$ (50-digit) & worst ratio $1.000000$\\
Prop.~\ref{prop:bounds} & $|\sigma_H^2/\Ain-1|\le\bar\Delta$ & worst ratio $1.000000$\\
Prop.~\ref{prop:bounds} & $c_H/\Ain\le\bar\gamma$ & worst ratio $1.000000$\\
Prop.~\ref{prop:bounds} & $|\sigma_{\dot H}^2/\Bin-1|\le(\lambda/\kappa)\bar\Delta$ & worst ratio $0.997226$\\
Prop.~\ref{prop:bounds} & tightened $\bar\Delta,\bar\gamma,\bar e$ (50-digit) & worst ratios $1.000000$ (all three)\\
Prop.~\ref{prop:point} & $|\ln\nu-\ln\nustat|\le\mathcal{E}$, $71\,079$ points & $0$ violations, worst ratio $0.612$\\
Prop.~\ref{prop:transfer} & the two functional norms $1.504$, $2.916$ & computed exactly\\
\bottomrule
\end{tabular}
\end{center}

Grids: $\lambda\in[0.3,60]$, $\kappa\in[0.01,60]$, $\rho\in[0.05,6]$,
$\beta/\delta\in\{0,0.05,0.3,1,3\}$, $a_\infty\in\{0.5,1,2,3.3,6,10\}$. The $X$ bounds
were verified in 50-digit arithmetic because the combination
$P_{\varepsilon\varepsilon}-2P_{\varepsilon H}+P_{HH}$ cancels catastrophically in
double precision; points where $\mathcal{E}$ falls below double-precision resolution
($10^{-10}$) were excluded from the Theorem~\ref{prop:point} check for the same reason.

\section{Consequences for the paper}

Adopting \eqref{eq:warm} removes from the appendix the convolution kernel $K$, the
integrals $\int K^2$ and $\int K'^2$, the auxiliary $W=K'-K^3/2\sigma_H^2$, the
log-concavity lemma, the non-stationary derivative of $\ln\nu$, and the transient
content of the present Corollaries A.2 and A.3. What remains is
Sections~\ref{sec:moments}--\ref{sec:neglect} --- five short results establishing that
the transient may be neglected --- followed by the stationary rate and the predictions
as one-line corollaries.

The cost is that theory and simulation must agree, so the $1200$-draw parameter sweep,
the admissibility screen and every effect size would have to be recomputed under
\eqref{eq:warm}.

\end{document}
